% Generated draft: requires mathematical and source review.
\documentclass[11pt]{article}
\usepackage[margin=25mm]{geometry}
\usepackage{fontspec}
\setmainfont{Latin Modern Roman}
\usepackage{amsmath,amssymb,mathtools,unicode-math}
\setmathfont{Latin Modern Math}
\usepackage{enumitem,needspace,xurl,hyperref}
\setlength{\parindent}{0pt}
\setlength{\parskip}{5pt}
\setlength{\emergencystretch}{3em}
\newcommand{\N}{\mathbb{N}}
\newcommand{\Z}{\mathbb{Z}}
\newcommand{\Q}{\mathbb{Q}}
\newcommand{\R}{\mathbb{R}}
\newcommand{\C}{\mathbb{C}}
\newcommand{\F}{\mathbb{F}}
\newcommand{\A}{\mathbb{A}}
\newcommand{\PP}{\mathbb{P}}
\newcommand{\E}{\mathbb{E}}
\newcommand{\eps}{\varepsilon}
\newcommand{\dd}{\,\mathrm{d}}
\begin{document}
\section*{UnsolvedMath ID: 30003683}
\textbf{Source problem code:} catalogue entry
\subsection{Definitions and mathematical statement}
The supplied record concerns monks, monasteries, and mathematical instruments as objects of historical investigation. Its title and statement-verification note place the inquiry in medieval monasteries, but it specifies neither chronological or geographical boundaries nor a definition of mathematical instrument or a class of monastic communities.

The recoverable historical question is: What motivated monks to study mathematical instruments, and how were those instruments studied and used in monasteries?

This is a question about historical evidence, motivations, and practices. The supplied record provides no mathematical objects, hypotheses, quantifiers, or conclusion from which a precise mathematical proposition can be recovered. Its statement status is \texttt{unrecoverable}, and no separate clean statement is supplied. Consequently, the question is retained as a historical research inquiry; a mathematical formulation would require information absent from the record.
\subsection{Short English statement}
The inquiry asks why monks studied mathematical instruments and how monasteries studied and used them, with a medieval context supplied by the title and verification note. No precise mathematical statement is recoverable from the supplied record.
\subsection{Sources}
\begin{itemize}
\item Local UnsolvedMath registry, supplied complete source record for permanent ID \(30003683\), Oberwolfach Reports Open Problems collection, source code \texttt{catalogue entry}. The original question is preserved; the clean statement is absent and the statement status is \texttt{unrecoverable}.
\item Supplied source citation: \emph{Mathematical Instruments between Material Artifacts and Ideal Machines: Their Scientific and Social Role before 1950} (2018), Oberwolfach Reports, pp.\ 3471--3560, DOI \texttt{10.4171/owr/2017/58}. Source URL and supplied literature-source link, labelled ``Oberwolfach report'': \url{https://doi.org/10.4171/owr/2017/58}. This report is supplied as provenance for workshop discussions; its contents have not been independently inspected here.
\item Attribution: no individual proposer is supplied. The registry gives the proposed year as 2017 and the report citation year as 2018; both are retained as supplied.
\item Supplied literature assessment, checked on exactly \texttt{2026-08-21}: the registry labels the inquiry ``open'' and states that there is no complete account in the literature of monastic motivations, study, and use of instruments. The listed source checked is the Oberwolfach report linked above. This is supplied dated curation derived from the OWR record, not an independently verified assessment of the literature or a verified mathematical open-problem status; time-sensitive status and cited claims require independent verification.
\end{itemize}
\end{document}
